\section{Task Examples}
\label{app:tasks}
\begin{figure}[t]
\centering\includegraphics[width=\linewidth]{figures/fig_g1_benchmark.png}
\caption{Representative task--output pairs. Requests are abridged and translated; source slide images retain their original language.}
\label{fig:tasks}
\end{figure}

\section{Detailed Criteria}
\label{app:taxonomy}
\paragraph{Instruction categories.}
Count and format cover numerical specifications such as slide count and aspect ratio. Language covers the requested output language. Layout and structure cover explicitly requested organization. Mandatory content concerns items explicitly required by the request, distinguished from broader source coverage. Template retention concerns visible compliance with a specified template or style. Additions and revisions concern edits to an existing deck, checked against count and content requirements. Missing evidence can make a criterion not applicable; NA items do not enter the pass-rate denominator.

\paragraph{Visual hard defects: four binary checks.}
Boundary cropping denotes content clipped by or extending beyond page boundaries. Occlusion/overlap denotes elements blocking one another and impairing recognition. Content integrity covers incomplete text or images, truncation, empty frames, or unreplaced placeholders. Image quality covers insufficient resolution, distortion, placeholder frames, or spurious images. Only cropping and occlusion enter the visual branch of the blocking-failure union; the other checks remain in visual reporting.

\paragraph{Fine-grained design: six 0/1/2 criteria.}
Overall alignment checks margins and grids. Spacing balance checks crowding and uneven gaps. Component consistency checks fonts, sizes, corners, and emphasis colors for comparable elements. Legibility checks text clarity and foreground--background contrast. Type hierarchy checks separation and consistency of title, subtitle, and body sizes. Leading and text flow check line spacing, overflow, wrapping, and broken text.

\paragraph{Holistic aesthetics: three 0--3 criteria.}
Layout hierarchy and reading path concern stable composition, emphasis, and visual order. Visual language concerns coherent colors, typography, and assets. Overall impression concerns apparent finish and completeness. Cross-page consistency is assessed at deck level. Normalized scores are combined using the formula in Section~4; per-type defect-page counts are retained separately.

\section{Initial Grouping and Screening Cohorts}
\label{app:groups}
The initial grouping contains 235 commonly scored decks: 124 with detected failures and 111 without detected failures. The latter group includes 29 with incomplete failure judgments. The abstract's five nonsignificant comparisons refer specifically to this initial grouping. Table~\ref{tab:initial} records the uncorrected two-sided Mann--Whitney results. Statistical nonsignificance does not demonstrate equivalence or the absence of a failure-related signal.

\begin{table}[h]
\caption{Initial grouping results (235 outputs). Unknown states are included in the no-detected-failure group in this analysis only.}
\label{tab:initial}\centering\small
\begin{tabular}{@{}lr@{}}
\toprule Scoring scheme & Uncorrected $p$\\\midrule
SlideAudit & .292\\
LecEval & .586\\
PPTBench-style & .001279\\
PPTEval & .061858\\
AutoPresent & .969\\
SlidesGen & .602\\\bottomrule
\end{tabular}
\end{table}

The primary screening cohort excludes the 29 unknown states, giving 206 decks with 124 positive and 82 negative labels. A positive label requires at least one detected blocking failure even if another dimension is unavailable; a negative requires evaluable judgments in all three dimensions and no detected failure. Requiring completeness in all dimensions also for positives yields the 186-deck sensitivity cohort (104 positive, 82 negative). All cohorts are subsets of the 240-output collection. These distinctions explain why the initial group tests and held-out screening are not analyses of identical samples.

For each held-out task, threshold selection uses only the remaining tasks. A deck is flagged when its quality score is at or below the selected cutoff. The cutoff is the strictest satisfying at least 90\% training recall; ties may produce higher recall. The held-out outcomes yield the confusion matrices in Table~\ref{tab:screening}. This operating point is an analysis choice rather than an official threshold of any compared framework.

\section{Ranking Conflicts}
\label{app:conflicts}
\begin{figure}[h]
\centering\includegraphics[width=.88\linewidth]{figures/fig_f2_conflicts.png}
\caption{Ranking-direction conflicts on the historical complete subset of 193 decks. The within-task permutation baseline is close to the observed rates.}
\end{figure}
The conflict analysis uses its own historical complete subset, not the 206-deck screening cohort. It reports opposing signs of score differences for system pairs on the same task. The result is descriptive and is not evidence that the dimensions depart significantly from an independence baseline.

\section{Aggregation Analyses}
\label{app:aggregation}
\paragraph{Visual penalty ablation.}
The existing visual ablation uses a separate scoring snapshot. Removing the linear penalty ($\lambda=0$) gives a system-score range of .103 and standard deviation .052. The current linear penalty ($\lambda=.8$) expands these to .141 and .071. Correlations with two Bradley--Terry references on that snapshot are .952 and .943. Multiplicative gating by .65 gives correlation .794 and standard deviation .109. These snapshot-specific correlations must not be substituted for the .855 validation result in the abstract.

\paragraph{Cross-dimension weights.}
A simplex grid with step .1 produces 66 weight settings on the historical 193-deck complete subset, yielding 22 distinct full rankings. Kimi ranks first in 40 settings (60.6\%), AnyGen in 21 (31.8\%), and Manus in five (7.6\%). The minimum Kendall correlation with equal-weight ranking is .067. These results concern weighting across dimensions; they do not establish sensitivity to every coefficient within a dimension.

\section{Source-Grounded Evidence}
\label{app:evidence}
\begin{figure}[h]
\centering\includegraphics[width=\linewidth]{figures/fig_g2_evidence.png}
\caption{Page-level source evidence in a training task. The two required activities can be located in a correctly grounded output, while the full Doubao record marks them missing. Original source and output images retain their original language.}
\end{figure}
The coverage illustration in Figure~\ref{fig:coverage} applies seven scored items to both decks. Manus receives six covered and one partial verdict, for $(6+0.5)/7=.929$ coverage. Doubao receives one partial and six missing verdicts, for $.5/7=.071$. For the item requiring both activities, the evidence is on Manus page 9 and absent across both evaluation batches of the Doubao deck. This example illustrates the distinction between a visible topic mention and coverage of required source content.
